一个风险评分系统上线后，两份分析报告同时说自己发现了真相：一份说在每个分数档上两个群体的实际结果比例一致，另一份说在最终没有出事的人里两个群体被判高风险的比例相差近一倍。两份报告读的是同一张表，算的都对，结论相反。

这不是谁算错了，是``公平''这个词在被两种方式定义。一种定义固定分数去看结果，另一种定义固定结果去看分数——条件写在竖线的哪一边，量就换了一个。而这类定义不止两个，每一个都对应一句听上去无可辩驳的正义直觉，指向的却是联合分布的不同边缘。

本章不打算在这些定义里评出高下，那不是数学能做的事。它做的是另一件更有用的事：把这些定义之间的关系算清楚。结果是一条硬性的限制——只要两个群体的基率不同、且分类器会犯错，校准、等真阳率、等假阳率就不能同时成立。这条限制来自一个四行的代数恒等式，不来自模型质量，因此加数据、换模型族、调超参都推不开它。能选的只有让哪一条不成立。

这条命题最常被读成``公平不可能''。它的实际内容要窄得多，也有用得多：``公平''必须先被指定成一个具体判据，指定之后才谈得上做到与否；不指定就争论，双方会永远各握一个真数字各说各话。

而所有这些判据还共享一个更深的限制：它们只依赖联合分布。人们真正想问的往往是``若这个人属于另一个群体，决定会不会变''，那是反事实层面的问题，观察分布回答不了。跨过去需要一张因果图，图不能从数据里识别，只能由领域知识声明——于是争论从判据之争变成图之争。这一步必须诚实说出来：问题没有消失，只是换了地方。

\subsection{三节怎么安排}

第一节把开头那场争论拆成两个方向相反的条件概率，并把主要判据并排摆开，看清它们在对同一张 \(2\times2\) 表的不同分母做归一化。第二节是本章的支点，给出不可能性命题的完整证明与条件对账，说明基率差与不完美这两个前提各自挡住什么误读。第三节走到统计判据的边界之外，指出所有观察层判据都区分不了``差异来自真实风险''与``差异来自直接进入决策的群体属性''，并交代反事实公平要求什么、又把难题推到了哪里。

\subsection{本章篇目}

\begin{itemize}
\tightlist
\item \hyperref[sec:17-Constrained-Guarantees/03-Fairness-Impossibility/01]{``同一个模型可以既公平又不公平''}；
\item \hyperref[sec:17-Constrained-Guarantees/03-Fairness-Impossibility/02]{``三个合理要求不能同时满足''}；
\item \hyperref[sec:17-Constrained-Guarantees/03-Fairness-Impossibility/03]{``统计公平到此为止''}。
\end{itemize}

读这一章需要的前置只有两样：条件概率两个方向不能互换，见\hyperref[sec:05-Probability/02-Conditional-Probability/02]{``全概率公式与 Bayes 更新''}；以及观察、干预、反事实分属三层，见\hyperref[ch:120]{``因果推断：预测世界不等于改变世界''}。第二节那个恒等式本身只用到 \(2\times2\) 表的边缘化，不需要任何额外工具。
